\documentclass[11pt]{article}

\usepackage{amsmath,amssymb,amsfonts}
\usepackage{physics}
\usepackage{geometry}
\usepackage{hyperref}
\usepackage{cite}
\usepackage{bm}
\title{Coordinate Eigenstates and Position Measurement}
\author{Shimon Panfil}
\date{}

\begin{document}
\maketitle

\section{Axiomatic Setting}

\begin{axiom}[Quantum kinematics]
A quantum system is described by a complex separable Hilbert space $\mathcal H$.
Physical observables are represented by normalized positive operator--valued measures (POVMs) on $\mathcal H$.
\end{axiom}

\begin{axiom}[Configuration space]
The configuration space of a nonrelativistic particle is a measurable space $(\Omega,\mathcal B(\Omega))$,
where $\Omega \subseteq \mathbb R^d$ may be unbounded or bounded.
\end{axiom}

\begin{axiom}[Position observable]
The position observable is represented by the POVM
\[
E_X(\Delta)\psi(x) = \chi_\Delta(x)\psi(x),
\quad \Delta \in \mathcal B(\Omega),
\]
acting on $\mathcal H = L^2(\Omega)$.
When $\Omega \subseteq \mathbb R^d$, this POVM is a projection-valued measure (PVM).
\end{axiom}

---

\section{The Position Operator}

\begin{definition}[Position operator]
For each coordinate $i=1,\dots,d$, define
\[
(\hat x^i \psi)(x) = x^i \psi(x)
\]
with maximal domain
\[
D(\hat x^i)=\{\psi\in L^2(\Omega)\mid x^i\psi(x)\in L^2(\Omega)\}.
\]
\end{definition}

\begin{proposition}
Each $\hat x^i$ is a self-adjoint multiplication operator with spectrum
\[
\sigma(\hat x^i)=\overline{\{x^i:x\in\Omega\}}.
\]
\end{proposition}

---

\section{Generalized Coordinate Eigenstates}

\begin{definition}[Coordinate eigenstate]
A coordinate eigenstate at $x_0\in\Omega$ is a distribution
$|x_0\rangle \in \Phi^\times$ such that
\[
\hat x^i |x_0\rangle = x_0^i |x_0\rangle
\quad \forall i.
\]
\end{definition}

\begin{theorem}[Existence and form]
For each $x_0\in\Omega$ there exists a generalized eigenstate $|x_0\rangle$ satisfying
\[
\langle x|x_0\rangle = \delta(x-x_0),
\]
where $\delta$ is the Dirac distribution restricted to $\Omega$.
\end{theorem}

\begin{proof}
Let $\Phi\subset L^2(\Omega)$ be a dense test space (e.g.\ $C_c^\infty(\Omega)$).
Define the functional
\[
\langle x_0|\phi\rangle := \phi(x_0).
\]
This functional is continuous on $\Phi$ and satisfies the eigenvalue equation in the distributional sense.
\end{proof}

\begin{corollary}
No coordinate eigenstate belongs to $\mathcal H$.
\end{corollary}

---

\section{Rigged Hilbert Space Structure}

\begin{definition}[Gelfand triple]
A rigged Hilbert space is a triple
\[
\Phi \subset \mathcal H \subset \Phi^\times
\]
where $\Phi$ is a dense nuclear space and $\Phi^\times$ its continuous dual.
\end{definition}

\begin{proposition}
Coordinate eigenstates $|x\rangle$ belong to $\Phi^\times$ and not to $\mathcal H$.
\end{proposition}

This construction formalizes the distributional calculus introduced by
\emph{:contentReference[oaicite:0]{index=0}}.

---

\section{Resolution of the Identity}

\begin{theorem}[Completeness]
The family $\{|x\rangle\}_{x\in\Omega}$ satisfies
\[
\int_\Omega |x\rangle\langle x|\,dx = \mathbb I
\]
in the weak operator sense.
\end{theorem}

\begin{corollary}
For any $\psi\in\mathcal H$,
\[
\psi(x) = \langle x|\psi\rangle .
\]
\end{corollary}

---

\section{Bounded Configuration Spaces}

\begin{theorem}[Position on bounded domains]
Let $\Omega\subset\mathbb R^d$ be bounded.
Then:
\begin{enumerate}
\item The position POVM remains a PVM.
\item The spectrum of $\hat x$ is $\Omega$.
\item Coordinate eigenstates exist as distributions.
\item No self-adjointness obstruction arises.
\end{enumerate}
\end{theorem}

\begin{remark}
This sharply contrasts with momentum observables, whose generators generally fail
to be self-adjoint on bounded domains.
\end{remark}

---

\section{Dual Structure: Momentum Eigenstates}

\begin{theorem}[Momentum on $\mathbb R^d$]
On $\mathcal H=L^2(\mathbb R^d)$, the operator
\[
\hat p=-i\hbar\nabla
\]
admits generalized eigenstates $|p\rangle$ with
\[
\langle x|p\rangle=(2\pi\hbar)^{-d/2}e^{ip\cdot x/\hbar}.
\]
\end{theorem}

\begin{remark}
On bounded $\Omega$, no such eigenstates exist in general; momentum measurements must
be described by POVMs rather than PVMs.
\end{remark}

---

\section{Measurement and Collapse}

\begin{axiom}[Operational measurement]
A position measurement with finite resolution $\Delta$ updates the state $\rho$ by
\[
\rho \mapsto \frac{E_X(\Delta)\rho E_X(\Delta)}{\mathrm{Tr}(\rho E_X(\Delta))}.
\]
\end{axiom}

\begin{theorem}[No exact collapse]
No physical measurement produces a coordinate eigenstate.
Exact localization corresponds to a singular limit outside $\mathcal H$.
\end{theorem}

\begin{proof}
$E_X(\Delta)$ maps $\mathcal H$ into itself for all measurable $\Delta$.
The limit $\Delta\to\{x_0\}$ is distributional and exits $\mathcal H$.
\end{proof}

---

\section{Physical Interpretation}

\begin{proposition}
Coordinate eigenstates do not represent physical states but idealized measurement outcomes.
\end{proposition}

\begin{remark}
Collapse is not a dynamical process in Hilbert space but an operational conditioning
arising from system--apparatus entanglement and decoherence.
\end{remark}

---

\section{Summary}

\begin{itemize}
\item Coordinate eigenstates are generalized vectors.
\item They define the position observable but are not realizable states.
\item Position remains well-defined on bounded domains.
\item Collapse to a point is a mathematical idealization.
\end{itemize}
\end{document}
